%% =====================================================================
%%  structures/informe-sr/componentes.tex  --  Componentes editoriales
%% ---------------------------------------------------------------------
%%  EDICION PUBLICA (revista): solo tipografia, aire y fotografia.
%%     \isrkicker{texto}                    etiqueta de seccion (sin regla)
%%     \isrcifra[ancho]{cifra}{etiqueta}{fuente}   dato grande (Inter)
%%     \isrcita{texto}{autor}{contexto}     cita destacada
%%     \isrfotoapertura{archivo}{pie}       foto grande de apertura (ancho de texto)
%%     \isrfotogrande{archivo}{pie}         foto a ancho de texto
%%     \isrfotosangre{archivo}{pie}         foto a sangre horizontal (21 cm)
%%     \begin{isrfotoytexto}{archivo}{pie} ... \end{isrfotoytexto}
%%                                          foto vertical + prosa al lado
%%     \isrfotocolumna{archivo}{pie}         foto al ancho de la columna, DENTRO
%%                                          de columnas (segunda foto de una cronica)
%%     \isrgaleriatitulo{texto}             kicker de la galeria del dia (pide sitio)
%%     \isrgaleriafila{a1}{pie1}{a2}{pie2}  fila indivisible de dos fotos 3:2 (a2 opcional)
%%     \isrgaleriatres{a1}{p1}{a2}{p2}{a3}{p3}  fila de tres (edicion a una columna)
%%     \isrfotoslot[ancho]{alto}{codigo}    espacio preparado, discreto
%%     \begin{columnas} ... \end{columnas}  prosa a dos columnas (una con \isrunacolumnatrue)
%%     \begin{isrpedidos} \item[P-011] ...  lista con codigo al margen
%%     \begin{isrcrono} \item[24 ago] ...   lista con fecha al margen
%%     \isrdato{etiqueta}{valor}            linea de ficha tecnica
%%     \isrgrafico[ancho]{archivo}{titulo}{fuente}
%%     \isrfuente{texto}
%%  EDICION INTERNA (documento de trabajo): fichas de actividad, cajas,
%%  chips (sin lineas verticales; pedido de Edison 2026-09-03).
%% =====================================================================

%% --- Kicker (sin regla) --------------------------------------------------
\newcommand{\isrkicker}[1]{%
  \par\needspace{4\baselineskip}\noindent
  {\isrtxkicker\color{colorisrprimario}\MakeUppercase{#1}}\par
  \vspace{5pt}}

%% --- Cifra (dato grande) ---------------------------------------------------
\newcommand{\isrcifra}[4][0.235\linewidth]{%
  \begin{minipage}[t]{#1}
    \setlength{\parskip}{0pt}\raggedright
    {\fuentedatossemi\fontsize{30}{32}\selectfont\color{colorisrprimario}#2\par}
    \vspace{4pt}
    {\fontsize{8.8}{11}\selectfont\color{colorisrtexto}#3\par}
    \if\relax\detokenize{#4}\relax\else
      \vspace{2pt}{\isrtxpie Fuente: #4\par}\fi
  \end{minipage}}

%% --- Linea de fuente ------------------------------------------------------
\newcommand{\isrfuente}[1]{\par\vspace{2pt}\noindent{\isrtxpie Fuente: #1\par}}

%% --- Cita destacada --------------------------------------------------------
\newcommand{\isrcita}[3]{%
  \par\needspace{6\baselineskip}\vspace{12pt}\noindent
  \begin{minipage}{\linewidth}
    \setlength{\parskip}{0pt}%
    {\fontsize{13.5}{17.5}\selectfont\itshape\color{colorisrprofundo}«#1»\par}
    \vspace{7pt}
    {\isrtxkicker\color{colorisrtextosec}\MakeUppercase{#2}%
      \if\relax\detokenize{#3}\relax\else\ \textperiodcentered\ #3\fi\par}
  \end{minipage}\par\vspace{12pt}}

%% --- Fotografias (sin marcos) ------------------------------------------------
\newcommand{\isrfotogrande}[2]{%
  \par\needspace{0.42\textheight}\vspace{8pt}\noindent
  \begin{minipage}{\linewidth}
    \setlength{\parskip}{0pt}%
    \includegraphics[width=\linewidth]{#1}\par\vspace{4pt}
    \if\relax\detokenize{#2}\relax\else{\isrtxpie #2\par}\fi
  \end{minipage}\par\vspace{12pt}}

%%  Apertura de cronica: exige casi la pagina completa libre; si no la
%%  hay, salta de pagina para que la foto abra arriba con su titular.
\newcommand{\isrfotoapertura}[2]{%
  \par\Needspace{0.85\textheight}\vspace{4pt}\noindent
  \begin{minipage}{\linewidth}
    \setlength{\parskip}{0pt}%
    \includegraphics[width=\linewidth]{#1}\par\vspace{4pt}
    \if\relax\detokenize{#2}\relax\else{\isrtxpie #2\par}\fi
  \end{minipage}\par}

%%  A sangre horizontal: la caja sobresale 2 cm por cada lado (margenes
%%  de 20 mm) y ocupa exactamente el ancho del papel. Va en el flujo.
\newcommand{\isrfotosangre}[2]{%
  \par\needspace{0.45\textheight}\vspace{8pt}\noindent
  \makebox[\linewidth][c]{\includegraphics[width=\paperwidth]{#1}}\par\vspace{4pt}
  \if\relax\detokenize{#2}\relax\else\noindent{\isrtxpie #2\par}\fi
  \vspace{12pt}}

%%  Foto vertical (2:3) a la izquierda + prosa a la derecha
\newenvironment{isrfotoytexto}[2]{%
  \par\needspace{0.5\textheight}\vspace{8pt}\noindent
  \begin{minipage}[t]{0.40\linewidth}\vspace{0pt}
    \setlength{\parskip}{0pt}%
    \includegraphics[width=\linewidth]{#1}\par\vspace{4pt}
    {\isrtxpie #2\par}
  \end{minipage}\hfill
  \begin{minipage}[t]{0.55\linewidth}\vspace{0pt}\isrtxcolumna\setlength{\parskip}{5pt}%
}{%
  \end{minipage}\par\vspace{12pt}}

%%  Foto de columna: va DENTRO de \begin{columnas}, al ancho de la columna,
%%  y fluye con la prosa (multicol la mueve a la columna siguiente si no
%%  cabe). Es la segunda foto de una cronica: la apertura sigue siendo la
%%  foto grande; con esta se admiten como maximo dos por cronica
%%  (pedido de Edison 2026-09-11). Sin marco, pie en \isrtxpie.
\newcommand{\isrfotocolumna}[2]{%
  \par\vspace{6pt}\noindent
  \begin{minipage}{\dimexpr\linewidth-\leftskip-\rightskip\relax}
    \setlength{\parskip}{0pt}%
    \includegraphics[width=\linewidth]{#1}\par\vspace{4pt}
    \if\relax\detokenize{#2}\relax\else{\isrtxpie #2\par}\fi
  \end{minipage}\par\vspace{10pt}}

%%  Galeria del dia (al menos una foto por actividad; pedido de Edison
%%  2026-09-11): filas de dos fotos al mismo formato 3:2, sin marco, con
%%  pie corto. Cada fila es una unidad indivisible que pide su espacio: si
%%  no cabe, pasa entera a la pagina siguiente. El segundo par puede ir
%%  vacio (numero impar de fotos). El titulo pide sitio para si y una fila,
%%  para no quedar huerfano al pie de pagina.
\newcommand{\isrgaleriatitulo}[1]{\par\Needspace{0.34\textheight}\vspace{10pt}\isrkicker{#1}}
\newcommand{\isrgaleriafila}[4]{%
  \par\Needspace{0.28\textheight}\vspace{4pt}\noindent
  \begin{minipage}[t]{0.49\linewidth}\vspace{0pt}
    \setlength{\parskip}{0pt}%
    \includegraphics[width=\linewidth]{#1}\par\vspace{3pt}
    \if\relax\detokenize{#2}\relax\else{\isrtxpie #2\par}\fi
  \end{minipage}\hfill
  \begin{minipage}[t]{0.49\linewidth}\vspace{0pt}
    \setlength{\parskip}{0pt}%
    \if\relax\detokenize{#3}\relax\else
      \includegraphics[width=\linewidth]{#3}\par\vspace{3pt}
      \if\relax\detokenize{#4}\relax\else{\isrtxpie #4\par}\fi
    \fi
  \end{minipage}\par\vspace{8pt}}

%%  Fila de tres: disponible, pero NO se usa en el ISR (Edison 2026-09-11: la
%%  galeria va siempre a dos columnas, nunca tres o mas; con dos fotos, grandes).
\newcommand{\isrgaleriatile}[3]{%
  \begin{minipage}[t]{#1}\vspace{0pt}
    \setlength{\parskip}{0pt}%
    \if\relax\detokenize{#2}\relax\else
      \includegraphics[width=\linewidth]{#2}\par\vspace{3pt}
      \if\relax\detokenize{#3}\relax\else{\isrtxpie #3\par}\fi
    \fi
  \end{minipage}}
\newcommand{\isrgaleriatres}[6]{%
  \par\Needspace{0.22\textheight}\vspace{4pt}\noindent
  \isrgaleriatile{0.315\linewidth}{#1}{#2}\hfill
  \isrgaleriatile{0.315\linewidth}{#3}{#4}\hfill
  \isrgaleriatile{0.315\linewidth}{#5}{#6}\par\vspace{8pt}}

%%  Espacio preparado (discreto): campo gris claro con etiqueta
\newcommand{\isrfotoslot}[3][\linewidth]{%
  \begin{minipage}[t]{#1}
    \begin{tcolorbox}[enhanced, colback=neutral200, boxrule=0pt, arc=0pt,
        height=#2, halign=center, valign=center, boxsep=0pt, left=4pt, right=4pt,
        before skip=0pt, after skip=0pt]
      {\fuentedatossemi\fontsize{6.4}{9}\selectfont\color{colorisrtextosec}%
        \addfontfeatures{LetterSpace=12}FOTOGRAFÍA PENDIENTE
        \if\relax\detokenize{#3}\relax\else\ \textperiodcentered\ #3\fi\par}
    \end{tcolorbox}
  \end{minipage}}

%% --- Columnas de prosa ----------------------------------------------------
%%  Dos columnas (revista) por defecto. Con \isrunacolumnatrue en el preambulo
%%  (edicion a una columna, pedido de Edison 2026-09-11) la prosa va a una sola
%%  columna con la medida reducida 1 cm por lado (≈ 90 caracteres) y la
%%  interlinea algo mayor; fotos de apertura y galerias siguen a ancho completo.
\newif\ifisrunacolumna
\newenvironment{columnasdos}{\begin{multicols}{2}\isrtxcolumna}{\end{multicols}}
\newenvironment{columnasuna}{\par\begingroup\isrtxcolumnaancha
  \setlength{\leftskip}{1cm}\setlength{\rightskip}{1cm}}{\par\endgroup}
\newenvironment{columnas}{\ifisrunacolumna\columnasuna\else\columnasdos\fi}%
  {\ifisrunacolumna\endcolumnasuna\else\endcolumnasdos\fi}

%% --- Listas con etiqueta al margen ------------------------------------------
\newlist{isrpedidos}{itemize}{1}
\setlist[isrpedidos]{label={}, leftmargin=1.45cm, labelwidth=1.15cm, labelsep=0.30cm,
  align=left, itemsep=4pt, topsep=4pt, parsep=0pt, labelindent=0pt,
  font=\fuentedatossemi\fontsize{7.6}{10}\selectfont\color{colorisrprimario}}
\newlist{isrcrono}{itemize}{1}
\setlist[isrcrono]{label={}, leftmargin=2.35cm, labelwidth=2.05cm, labelsep=0.30cm,
  align=left, itemsep=4pt, topsep=4pt, parsep=0pt, labelindent=0pt,
  font=\fuentedatossemi\fontsize{7.6}{10}\selectfont\color{colorisrprimario}}

%% --- Linea de ficha tecnica (etiqueta pequena + valor) ----------------------
\newcommand{\isrdato}[2]{%
  \par\noindent{\isrtxkicker\color{colorisrtextosec}\MakeUppercase{#1}\par}%
  \vspace{1.5pt}%
  {\fontsize{9.2}{12.5}\selectfont\color{colorisrtexto}#2\par}%
  \vspace{7pt}}

%% --- Grafico -------------------------------------------------------------------
\newcommand{\isrgrafico}[4][\linewidth]{%
  \par\needspace{10\baselineskip}\vspace{8pt}\noindent
  \begin{minipage}{#1}
    \setlength{\parskip}{0pt}%
    \isrkicker{#3}%
    \includegraphics[width=\linewidth]{#2}\par\vspace{3pt}
    {\isrtxpie Fuente: #4\par}
  \end{minipage}\par\vspace{12pt}}

%% =====================================================================
%%  EDICION INTERNA: chips, cajas y fichas de actividad
%% =====================================================================
\newcommand{\isrchip}[2][colorisrprimario]{%
  {\setlength{\fboxsep}{2.4pt}\colorbox{#1}{%
    \fuentedatossemi\fontsize{6.2}{6.2}\selectfont\color{white}%
    \addfontfeatures{LetterSpace=4}\,#2\,}}}
\newcommand{\isrchipclaro}[2][colorisrfila]{%
  {\setlength{\fboxsep}{2.4pt}\colorbox{#1}{%
    \fuentedatossemi\fontsize{6.2}{6.2}\selectfont\color{colorisrtexto}%
    \addfontfeatures{LetterSpace=4}\,#2\,}}}
\newcommand{\isrprioridad}[1]{%
  \ifcsname isr-prio-#1\endcsname\csname isr-prio-#1\endcsname\else\isrchipclaro{#1}\fi}
\expandafter\def\csname isr-prio-urgente\endcsname{\isrchip[puka700]{URGENTE}}
\expandafter\def\csname isr-prio-alta\endcsname{\isrchip[qillu600]{ALTA}}
\expandafter\def\csname isr-prio-ordinaria\endcsname{\isrchipclaro[neutral300]{ORDINARIA}}
\newcommand{\isrestado}[1]{%
  \ifcsname isr-est-#1\endcsname\csname isr-est-#1\endcsname\else\isrchipclaro{#1}\fi}
\expandafter\def\csname isr-est-registrado\endcsname{\isrchipclaro[neutral200]{REGISTRADO}}
\expandafter\def\csname isr-est-derivado\endcsname{\isrchipclaro[qillu100]{DERIVADO}}
\expandafter\def\csname isr-est-respondido\endcsname{\isrchipclaro[puka100]{RESPONDIDO}}
\expandafter\def\csname isr-est-cerrado\endcsname{\isrchip[exito]{CERRADO}}

\definecolor{qillu700}{HTML}{83600D}
\tcbset{
  isrcajabase/.style={enhanced, breakable, arc=0pt, outer arc=0pt,
    boxrule=0pt, left=10pt, right=10pt, top=8pt, bottom=8pt, boxsep=0pt,
    before skip=9pt, after skip=9pt,
    fontupper=\isrtxcuerpo, coltext=colorisrtexto,
    attach title to upper=\par\vspace{4pt},
    fonttitle=\isrtxkicker, coltitle=colorisrprimario},
  isrcaja-nota/.style={colback=colorisrfondo},
  isrcaja-alerta/.style={colback=colorisrfondosuave},
  isrcaja-compromiso/.style={colback=qillu50, coltitle=qillu700},
  isrcaja-dato/.style={colback=white, colframe=colorisrborde, boxrule=0.5pt},
  isrcaja-oscura/.style={colback=colorisrprofundo, coltext=white, coltitle=puka200},
}
\newtcolorbox{isrcaja}[2][nota]{isrcajabase, isrcaja-#1, title={\MakeUppercase{#2}}}

\newcommand{\isrfoto}[4][\linewidth]{%
  \begin{minipage}[t]{#1}
    \setlength{\parskip}{0pt}%
    \includegraphics[width=\linewidth]{#2}\par
    \vspace{3pt}%
    \if\relax\detokenize{#3}\relax\else
      {\fontsize{8.2}{10}\selectfont\color{colorisrtexto}#3\par}\fi
    \if\relax\detokenize{#4}\relax\else
      {\isrtxpie #4\par}\fi
  \end{minipage}}

\pgfkeys{/isr/ficha/.is family, /isr/ficha,
  id/.store in=\isrFichaId,
  titulo/.store in=\isrFichaTitulo,
  tipo/.store in=\isrFichaTipo,
  sector/.store in=\isrFichaSector,
  hora/.store in=\isrFichaHora,
  duracion/.store in=\isrFichaDuracion,
  lugar/.store in=\isrFichaLugar,
  ubicacion/.store in=\isrFichaUbicacion,
  participantes/.store in=\isrFichaParticipantes,
  resumenpedidos/.store in=\isrFichaResumenPedidos,
  documentos/.store in=\isrFichaDocumentos,
  foto/.store in=\isrFichaFoto,
  fotopie/.store in=\isrFichaFotoPie,
  fotoancho/.store in=\isrFichaFotoAncho,
  observacion/.store in=\isrFichaObs,
  default/.style={id=, titulo=, tipo=, sector=gobierno, hora=, duracion=,
    lugar=, ubicacion=, participantes=, resumenpedidos=, documentos=,
    foto=, fotopie=, fotoancho=0.34, observacion=}}
\tcbset{isrficha/.style={enhanced, breakable=false, colback=white,
  colframe=colorisrborde, boxrule=0.4pt, arc=0pt, outer arc=0pt,
  left=10pt, right=10pt, top=7pt, bottom=7pt, boxsep=0pt,
  before skip=9pt, after skip=9pt}}
\newcommand{\isrmeta}[2]{%
  \raisebox{-0.065cm}{\isricono[tinta]{#1}{0.29cm}}\hspace{2.5pt}#2\hspace{9pt}}
\newcommand{\isrfichacabecera}{%
  \noindent
  \isrchip{\isrFichaId}%
  \hspace{6pt}\raisebox{-0.1cm}{\isrsectoricono{\isrFichaSector}{0.36cm}}\hspace{3pt}%
  {\isrtxkicker\color{colorisrtextosec}\MakeUppercase{\isrsectornombre{\isrFichaSector}}}%
  \hfill
  {\isrtxkicker\color{colorisrtextosec}\MakeUppercase{\isrFichaTipo}}%
  \par\vspace{5pt}
  {\fuentedisplaysemi\fontsize{11.5}{13.5}\selectfont\color{colorisrtexto}\isrFichaTitulo\par}
  \vspace{3.5pt}
  {\isrtxmeta
    \if\relax\detokenize\expandafter{\isrFichaHora}\relax\else
      \isrmeta{clock}{\isrFichaHora
        \if\relax\detokenize\expandafter{\isrFichaDuracion}\relax\else\ (\isrFichaDuracion)\fi}\fi
    \if\relax\detokenize\expandafter{\isrFichaLugar}\relax\else
      \isrmeta{map-pin}{\isrFichaLugar
        \if\relax\detokenize\expandafter{\isrFichaUbicacion}\relax\else
          \ \textperiodcentered\ \isrFichaUbicacion\fi}\fi
    \if\relax\detokenize\expandafter{\isrFichaParticipantes}\relax\else
      \isrmeta{users}{\isrFichaParticipantes}\fi
    \par}
  \vspace{7pt}}
\newcommand{\isrfichapie}{%
  \par\vspace{6pt}
  \noindent{\isrtxmeta
    \if\relax\detokenize\expandafter{\isrFichaResumenPedidos}\relax\else
      \isrFichaResumenPedidos\fi
    \if\relax\detokenize\expandafter{\isrFichaDocumentos}\relax\else
      \hspace{10pt}\isrFichaDocumentos\fi
    \par}%
  \if\relax\detokenize\expandafter{\isrFichaObs}\relax\else
    \vspace{2.5pt}\noindent{\fontsize{7.6}{9.5}\selectfont\itshape\color{colorisrtextosec}%
      Nota de la fuente: \isrFichaObs\par}\fi}
\def\isrFotoSlot{slot}
\newenvironment{fichaactividad}[1]{%
  \pgfkeys{/isr/ficha,default,#1}%
  \begin{tcolorbox}[isrficha]
    \isrfichacabecera
    \if\relax\detokenize\expandafter{\isrFichaFoto}\relax
      \def\isrFichaTextoAncho{\linewidth}%
    \else
      \def\isrFichaTextoAncho{\dimexpr\linewidth-\isrFichaFotoAncho\linewidth-12pt\relax}%
    \fi
    \noindent\begin{minipage}[t]{\isrFichaTextoAncho}
      \setlength{\parskip}{3.5pt}\isrtxcuerpo\RaggedRight
}{%
    \end{minipage}%
    \if\relax\detokenize\expandafter{\isrFichaFoto}\relax\else
      \hfill
      \begin{minipage}[t]{\isrFichaFotoAncho\linewidth}
        \vspace{0pt}%
        \ifx\isrFichaFoto\isrFotoSlot
          \isrfotoslot{3.4cm}{\isrFichaId}%
        \else
          \isrfoto{\isrFichaFoto}{\isrFichaFotoPie}{}%
        \fi
      \end{minipage}%
    \fi
    \isrfichapie
  \end{tcolorbox}}
\newcommand{\isrcontexto}[1]{\noindent#1\par}
\newcommand{\isrsubtitulo}[1]{%
  \par\vspace{4pt}{\isrtxkicker\color{colorisrprimario}\MakeUppercase{#1}\par}\vspace{1.5pt}}
\newlist{isrhallazgos}{itemize}{1}
\setlist[isrhallazgos]{
  label={\textcolor{colorisrprimario}{\rule[0.09em]{0.40em}{0.40em}}},
  leftmargin=1.05em, labelsep=0.5em, labelindent=0pt,
  itemsep=1.5pt, topsep=2pt, parsep=0pt}
